\documentclass[12pt]{article}
\usepackage[utf8]{inputenc}
\usepackage[T1]{fontenc}
\usepackage{hyperref}
\title{Yoy Language Resources for K-12 Curriculum Development}
\author{Compiled by Claude Sonnet 5.5 (claude-sonnet-5-5) for Andr\'as Kornai}
\date{}
\begin{document}
\maketitle

\section{Introduction}
Yoy (ISO 639-3 \texttt{yoy}; Glottocode yoyy1238) is an endangered Tai (Kra-Dai) language of northeastern Thailand and central Laos. Wikipedia cites about 6,000 speakers. Around 7,000 Yoy people live mainly in Sakon Nakhon province, and about 1,000 live in Khammouan province, Laos. The language is endangered because of rapid shift to Lao. All Yoy are reported to be bilingual in Lao. Not to be confused with Tai Yo, which is a different language.

\section{Primary Educational Resources}
No primers, textbooks, or school materials in Yoy were found in search results.

\section{Linguistic and Reference Materials}
\begin{itemize}
\item Wikipedia, ``Yoy language'': \url{https://en.wikipedia.org/wiki/Yoy_language}
\item Numeral data for Yoy from the Max Planck Institute for Evolutionary Anthropology (EVA) numerals database: \url{https://lingweb.eva.mpg.de/channumerals/Yoy.htm}
\end{itemize}

\section{Religious and Community Resources}
\begin{itemize}
\item Joshua Project, Yoy in Thailand: \url{https://joshuaproject.net/index.php/people_groups/16062/TH}
\item Joshua Project, Yoy in Laos: \url{https://joshuaproject.net/index.php/people_groups/16062/LA}
\item People-group profile (PDF): \url{https://people-groups.asiaharvest.org/Laos/Yoy-PBW.pdf}
\end{itemize}

\section{Digital and Audio Resources}
None found.

\section{Curriculum Development Context}
Few researchers have studied Yoy, and no pedagogical material was found. A curriculum would first need documentation work: an orthography (the neighbouring Lao and Thai scripts are the likely basis), a lexicon, and recordings. Nothing found indicates that the language is taught in schools.
\end{document}